% !TeX root = main.tex
% Synchronized with 00_摘要.md v0.6; Tsai format.
\begin{abstracten}
\begin{spacing}{1.5}

Mobile manipulation in dynamic environments requires obstacle avoidance, coordinated base and arm motion, and prediction that accounts for the difference between commanded and realized motion. This thesis develops a control architecture based on geometric prediction and execution dynamics modeling. Navigation uses an \(SE(2)\) pose-error model, predictive control barrier conditions, and scan-based protection. For coordinated manipulation, arm position setpoints and tracking dynamics are incorporated into prediction. Base and arm inputs are jointly optimized over a finite horizon to coordinate end-effector tracking and motion constraints, with visual measurements providing manipulation references.

In simulations over 83 routes evaluated across three methods, the proposed navigation method reached every goal, with contact on one route. MPPI and RPP reached 64 and 60 goals, respectively. In drawer manipulation, the coordinated strategy reduced task time by 13.38--14.05\% relative to a strict stationary-base strategy while satisfying the prescribed task criteria. Five-step prediction also yielded lower end-effector position errors over the evaluated control period and fewer arm-command reversals than one-step prediction. These results demonstrate that the proposed predictive and coordination designs can support manipulation quality and task efficiency under the tested conditions.

\keyworden{mobile manipulators, geometric model predictive control, whole-body coordination, execution dynamics modeling, control barrier functions, manipulation on the move}
\end{spacing}
\end{abstracten}
